% SPDX-License-Identifier: Apache-2.0
\section{Conclusion}
\label{sec:e-conclusion}

A companion paper merged two hardware description languages and gave the
result a syntax shaped after Rust. This paper removed every construct Rust
already provides and asked what was left.

Seventeen of twenty constructs go. What remains is two Rust modules of
traits and types, not a grammar.

Two deletions are worth more than the count suggests. The rule that a
timeless body may not count cycles needs no enforcement, because a cycle
count is spelled \code{.await} and a plain function has none: a paragraph
of specification and a compiler check become a fact about the type system.
And a pipeline, a sequence and a set of parallel processes turn out to be
one construct driven three ways, so \code{pipeline}, \code{stage} and
\code{pipe} all go, and the live-range analysis the last of them existed to
feed is what an async state machine does already.

Two results were negative and were worth more still, because neither was
visible from the language reference. Arbitrary-width arithmetic does not
work on stable Rust. A unit's parallel processes cannot take a mutable
borrow of the unit, and the repair, a register being a cell that several
processes reach, describes a register better than the version that failed.

One wall stands. A signal has no value while the program runs, so it cannot
be the condition of an \code{if}. That is what a separate grammar was
buying, and a macro front end over Rust's own syntax tree is the way to buy
it back.

The project's second thesis declined this direction on the grounds that a
full toolchain had become affordable. That argument has not weakened. It
now points the other way: if a toolchain is affordable, build the one that
arrives with a type system, a module system, generics and a package manager
already written.
